\scanpage{39}
\mspage{37}

Je vais \unsure{faire} quelques rappels catégoriques, généralisant la situation du foncteur canonique de localisation

\diagram{\begin{tikzcd}[ampersand replacement=\&]
 \mathrm{Cat\,hot}\arrow[r,"\varphi"'] \& \deleted{\unsure{C}}\mathrm{Hot},
\end{tikzcd}}

dont on sait qu'il admet un adjoint à droite

\diagram{\begin{tikzcd}[ampersand replacement=\&]
 \mathrm{Cat\,hot} \& \mathrm{Hot}\arrow[l,"\psi"'] \&[-1em] ,\quad
 \mathrm{Hom}_{\mathrm{Hot}}(\varphi X,Y)\simeq\mathrm{Hom}_{\mathrm{Cat\,hot}}(X,\psi Y)
\end{tikzcd}}

où $\psi$ est le foncteur ``enveloppe de Kan'' ou ``représentant de Kan''. On va donc considérer la situation plus générale

\diagram{\begin{tikzcd}[ampersand replacement=\&]
 C\arrow[r,shift right,"\varphi"'] \& C'\arrow[l,shift right,"\psi"'] \&[-1em]
 \mathrm{Hom}_{C'}(\varphi X,X')\simeq\mathrm{Hom}_C(X,\psi X')
\end{tikzcd}}

i.e. $\varphi,\psi$ couple de foncteurs adjoints.

On suppose de plus

\[
 \left\{\begin{aligned}
 &\varphi\ \text{un foncteur localisation, i.e. }\psi\ \text{pleinement fidèle},\\
 &\text{i.e.}\qquad\deleted{\unsure{\mathrm{id}_{C'}\longrightarrow}}
  \varphi\psi\xrightarrow{\sim}\mathrm{id}_{C'} .
 \end{aligned}\right.
\]

Soit $W=W_{\varphi}\subset\mathrm{Fl}\,C$, \unsure{syst} défini par $\varphi$: $W_\varphi=\{u\in\mathrm{Fl}\,C\mid\varphi(u)\ \mathrm{iso}\}$. Voici la proposition que j'avais oubliée, et que j'ai eu du mal à réconstituer, hélas !

\uline{Proposition 1} Soit $K\in\mathrm{Ob}\,C$. Les conditions suivantes \inserted{sur $K$} sont équivalentes :

(i) Le foncteur

\[
 X\longrightarrow K(X)\overset{\mathrm{def}}{=}\mathrm{Hom}_C(X,K)
 \qquad C^{\unsure{o|0}}\unsure{\relbar|\longrightarrow}\mathrm{Ens}
\]

transforme flèches de $W$ (``quasi-isomorphismes) en bijections

(i$'$) \unsure{Même|Mêmes} \unsure{condition|conditions} pour les flèches d'adjonction

\[
 X\longrightarrow\psi\varphi X\qquad(\forall X\in\mathrm{Ob}\,C)
\]

(i$\unsure{''}$) Même condition pour la flèche d'adjonction

\[
 K\longrightarrow\psi\varphi K
\]

\margin{N.B. Il suffit $=$ que l'élém $\mathrm{id}_K$ de $F_K(K)$ provienne de $\unsure{F_K}(\psi\varphi K)$ i.e. (i$'''$) Il existe une rétraction de $\psi\varphi(X)$ sur $X$ (\uline{NB} Celle-ci sera \unsure{même|bien} un iso)}

(ii) $K$ appartient à l'image essentielle de $\psi$

(ii$'$) $K\longrightarrow\psi\varphi K$ (adjonction) est un iso.

(ii$''$) L'hom fonctoriel en $X$

\[
 \mathrm{Hom}_C(X,K)\longrightarrow\mathrm{Hom}_{C'}(\varphi X,\varphi K)
\]

est un iso.

\uline{Dém} On a d'abord comme quasi-tautologies

\diagram{\begin{tikzcd}[ampersand replacement=\&,column sep=small,row sep=small]
 (\mathrm{ii}\unsure{'})\arrow[r,Rightarrow] \& (\mathrm{ii}\unsure{\vphantom{' }|'})\arrow[r,Rightarrow] \&(\mathrm{ii}'')\arrow[d,Rightarrow]\\
 \& \&(\mathrm{ii}')
\end{tikzcd}}

i.e. l'équivalence des trois dernières conditions. En effet (ii$'$) $\Longrightarrow$ (ii) tautologique, (ii) $\Longrightarrow$ (ii$''$) car si $K=\psi(K')$ on a

% archive sheet number: 38
